%------------------------------------------------
\section{Нагрузка током}
%------------------------------------------------

\begin{frame}
	\frametitle{Нагрузка током в линейном резонансном ускорителе электронов}
	\par Эффект нагрузки током состоит в том, что с повышением тока 
	ускоряемого пучка растет отбираемая пучком доля запасенной в 
	ускоряющей структуре энергии, что может приводить к значительному 
	снижению амплитуды ускоряющего поля.
	\newline
	\par  Для ускорителей электронов нагрузка током существенна, а влияние 
	кулоновского поля мало из-за быстрого достижения пучком 
	квазирелятивистских скоростей.
	\par Увеличивается и величина 
	высокочастотного поля, наведенного частицами в резонансной 
	ускоряющей системе, причем генерация собственного поля происходит во 
	всех полосах пропускания резонансной системы.
	\newline
	\par Наведенное пучком в ускоряющей структуре высокочастотное поле в общем случае зависит от 
	скорости частиц, формы и длительности токового импульса, то есть от 
	динамики частиц пучка.

\end{frame}

%------------------------------------------------

